\subsection{From Positional Access to Hash-Based Access}

The previous chapter studied Python objects whose components are arranged as sequences: \texttt{str}, \texttt{list}, and \texttt{tuple}. Their access model is positional. If an object is a sequence, the question ``where is the component?'' can be answered with an integer index.

Hash-based collections use a different access model. A \texttt{set} does not ask for the third element. A \texttt{dict} does not normally ask for the value at position three. Instead, these objects use hash-derived table placement. The stored object, or the stored key, is processed into an integer hash value, and that value guides Python toward a table slot where the relevant entry may be found.

This is the main transition of the chapter:

\begin{quote}
Sequential collections locate components by position. Hash-based collections locate components by value identity rules: hash first, equality confirmation after that.
\end{quote}

The difference is not cosmetic syntax. It changes what the structure is good at. A list is natural when order and position matter. A set is natural when uniqueness and membership matter. A dictionary is natural when a key must lead directly to an associated value.

\subsubsection{Sequential Collections: Why Lists, Tuples, and Strings Locate Components by Index}

A sequence is built around component order. The object has a first component, a second component, a third component, and so on. Python exposes that model through indexing, slicing, iteration, and length calculation.

\begin{verbatim}
text = "D'oh!"
items = ["spam", "eggs", 42, "larch"]
point = ("x", "y", "z")

print(f"text:  {text!r}")
print(f"items: {items!r}")
print(f"point: {point!r}")

print()
print(f"type(text):  {type(text)}")
print(f"type(items): {type(items)}")
print(f"type(point): {type(point)}")
\end{verbatim}

Possible output:

\begin{verbatim}
text:  "D'oh!"
items: ['spam', 'eggs', 42, 'larch']
point: ('x', 'y', 'z')

type(text):  <class 'str'>
type(items): <class 'list'>
type(point): <class 'tuple'>
\end{verbatim}

The common mechanism is visible if we inspect a few selected names from \texttt{dir()}.

\begin{verbatim}
text = "D'oh!"
items = ["spam", "eggs", 42, "larch"]
point = ("x", "y", "z")

names = ["__len__", "__getitem__", "__contains__", "index", "count"]

for obj_name, obj in [("text", text), ("items", items), ("point", point)]:
    print(f"{obj_name}: {type(obj)}")
    for name in names:
        print(f"  {name:<14} {name in dir(obj)}")
    print()
\end{verbatim}

Possible output:

\begin{verbatim}
text: <class 'str'>
  __len__        True
  __getitem__    True
  __contains__   True
  index          True
  count          True

items: <class 'list'>
  __len__        True
  __getitem__    True
  __contains__   True
  index          True
  count          True

point: <class 'tuple'>
  __len__        True
  __getitem__    True
  __contains__   True
  index          True
  count          True
\end{verbatim}

The important name here is \texttt{\_\_getitem\_\_}. It is the operation behind square-bracket access.

\begin{verbatim}
items = ["spam", "eggs", 42, "larch"]

print(f"items[0]:  {items[0]!r}")
print(f"items[1]:  {items[1]!r}")
print(f"items[-1]: {items[-1]!r}")
print(f"items[1:3]: {items[1:3]!r}")
\end{verbatim}

Possible output:

\begin{verbatim}
items[0]:  'spam'
items[1]:  'eggs'
items[-1]: 'larch'
items[1:3]: ['eggs', 42]
\end{verbatim}

For a sequence, the integer index is meaningful because the container has a positional structure. Index \texttt{0} means the first component. Index \texttt{1} means the second component. A slice such as \texttt{items[1:3]} means a new sequence containing the components found from one positional boundary to another.

Membership testing is available too, but it is not the same kind of access.

\begin{verbatim}
items = ["spam", "eggs", 42, "larch"]

print(f"{42 in items = }")
print(f"{'beans' in items = }")
print(f"{items.index(42) = }")
\end{verbatim}

Possible output:

\begin{verbatim}
42 in items = True
'beans' in items = False
items.index(42) = 2
\end{verbatim}

In a list, membership testing checks whether a value is present. The method \texttt{index()} then reports where that value was found. This still belongs to the positional world: after the value is discovered, Python can report a numeric location.

At the conceptual level, a sequence resembles a row of slots:

\begin{verbatim}
index:    0       1       2        3
value:  "spam"  "eggs"   42     "larch"
\end{verbatim}

The position is part of the structure. If two elements are swapped, the sequence has changed even if it still contains the same values.

\begin{verbatim}
a = ["spam", "eggs", 42]
b = [42, "eggs", "spam"]

print(f"a: {a!r}")
print(f"b: {b!r}")
print(f"a == b: {a == b}")
\end{verbatim}

Possible output:

\begin{verbatim}
a: ['spam', 'eggs', 42]
b: [42, 'eggs', 'spam']
a == b: False
\end{verbatim}

The two lists contain the same three values, but not in the same positions. For a sequence, that is a different structure.

\subsubsection{Hash-Based Collections: Why Sets and Dictionaries Locate Components by Hash-Derived Table Slots}

A hash-based collection is not organized around integer positions. It is organized around hashable values.

Python's main built-in hash-based collections are:

\begin{itemize}
    \item \texttt{set}: stores unique elements;
    \item \texttt{dict}: stores key-value associations.
\end{itemize}

A first inspection shows that these objects do not expose the ordinary sequence indexing model.

\begin{verbatim}
names = {"Arthur", "Brian", "Lancelot"}
quotes = {
    "Homer": "D'oh!",
    "Jerry": "What's the deal with airline food?",
    "Arthur": "It is only a model.",
}

print(f"names:  {names!r}")
print(f"quotes: {quotes!r}")

print()
print(f"type(names):  {type(names)}")
print(f"type(quotes): {type(quotes)}")
\end{verbatim}

Possible output:

\begin{verbatim}
names:  {'Brian', 'Arthur', 'Lancelot'}
quotes: {'Homer': "D'oh!", 'Jerry': "What's the deal with airline food?", 'Arthur': 'It is only a model.'}

type(names):  <class 'set'>
type(quotes): <class 'dict'>
\end{verbatim}

The displayed order of the set must not be read as the structural order of the set. A set is not a row of positions. It is a uniqueness domain: each accepted element is present once.

The selected \texttt{dir()} inspection also shows the change in available behavior.

\begin{verbatim}
names = {"Arthur", "Brian", "Lancelot"}
quotes = {
    "Homer": "D'oh!",
    "Jerry": "What's the deal with airline food?",
    "Arthur": "It is only a model.",
}

set_names = ["__len__", "__contains__", "__getitem__", "add", "remove", "union"]
dict_names = ["__len__", "__contains__", "__getitem__", "get", "keys", "items"]

print(f"names: {type(names)}")
for name in set_names:
    print(f"  {name:<14} {name in dir(names)}")

print()
print(f"quotes: {type(quotes)}")
for name in dict_names:
    print(f"  {name:<14} {name in dir(quotes)}")
\end{verbatim}

Possible output:

\begin{verbatim}
names: <class 'set'>
  __len__        True
  __contains__   True
  __getitem__    False
  add            True
  remove         True
  union          True

quotes: <class 'dict'>
  __len__        True
  __contains__   True
  __getitem__    True
  get            True
  keys           True
  items          True
\end{verbatim}

A \texttt{set} has no \texttt{\_\_getitem\_\_} operation because there is no meaningful \texttt{names[0]} access. The set can answer whether an element is present, but it is not designed to answer which element occupies the first, second, or third position.

\begin{verbatim}
names = {"Arthur", "Brian", "Lancelot"}

print(f"{'Arthur' in names = }")
print(f"{'Galahad' in names = }")

print()
try:
    print(names[0])
except TypeError as err:
    print(f"set indexing failed: {err}")
\end{verbatim}

Possible output:

\begin{verbatim}
'Arthur' in names = True
'Galahad' in names = False

set indexing failed: 'set' object is not subscriptable
\end{verbatim}

A dictionary does have \texttt{\_\_getitem\_\_}, but the square brackets do not mean positional access. In a dictionary, the bracket expression supplies a key, not an index.

\begin{verbatim}
quotes = {
    "Homer": "D'oh!",
    "Jerry": "What's the deal with airline food?",
    "Arthur": "It is only a model.",
}

print(f"quotes['Homer']:  {quotes['Homer']!r}")
print(f"quotes['Arthur']: {quotes['Arthur']!r}")

print()
print(f"{'Homer' in quotes = }")
print(f"{'D\\'oh!' in quotes = }")
\end{verbatim}

Possible output:

\begin{verbatim}
quotes['Homer']:  "D'oh!"
quotes['Arthur']: 'It is only a model.'

'Homer' in quotes = True
"D'oh!" in quotes = False
\end{verbatim}

The expression \texttt{'Homer' in quotes} checks whether \texttt{'Homer'} is a key. It does not check whether \texttt{'Homer'} appears somewhere among the values. This is one of the most important dictionary boundaries: membership belongs to the key domain.

The bridge from value to table slot is the hash value. Python exposes this through the built-in \texttt{hash()} function.

\begin{verbatim}
values = ["Arthur", "Brian", "Lancelot", 42]

for value in values:
    print(f"value={value!r:<12} type={type(value)} hash={hash(value)}")
\end{verbatim}

Possible output, with hash integers shortened here because some hash values may differ between Python runs:

\begin{verbatim}
value='Arthur'     type=<class 'str'> hash=...
value='Brian'      type=<class 'str'> hash=...
value='Lancelot'   type=<class 'str'> hash=...
value=42           type=<class 'int'> hash=42
\end{verbatim}

The hash value is not the stored value itself. It is an integer descriptor used by the hash table machinery to narrow the search area. For a set, the element is hashed. For a dictionary, the key is hashed.

Conceptually:

\begin{verbatim}
set lookup:
    element -> hash(element) -> candidate table slot -> equality check

dict lookup:
    key -> hash(key) -> candidate table slot -> equality check -> associated value
\end{verbatim}

This explains why a set can check membership without scanning from the first element to the last in the ordinary list sense. It also explains why a dictionary can retrieve a value by key without asking where that pair appears visually when printed.

\begin{verbatim}
names = {"Arthur", "Brian", "Lancelot"}
quotes = {
    "Homer": "D'oh!",
    "Jerry": "What's the deal with airline food?",
    "Arthur": "It is only a model.",
}

print("set membership")
print(f"{'Brian' in names = }")

print()
print("dictionary key lookup")
key = "Jerry"
print(f"key: {key!r}")
print(f"value: {quotes[key]!r}")
\end{verbatim}

Possible output:

\begin{verbatim}
set membership
'Brian' in names = True

dictionary key lookup
key: 'Jerry'
value: "What's the deal with airline food?"
\end{verbatim}

The practical distinction is therefore simple:

\begin{itemize}
    \item use a sequence when position is part of the meaning;
    \item use a set when uniqueness and membership are the main operations;
    \item use a dictionary when a key must lead to a value.
\end{itemize}

A list answers: ``what object is stored at this numeric position?''

A set answers: ``is this hashable object present?''

A dictionary answers: ``which value is associated with this hashable key?''